\begin{longtable}{lp{10cm}}
\toprule 符号 & 含义 \\ \midrule
$\T$ & 药材温度，内部K \\
$\C$ & 药材干基含水率，kg水/kg干物质 \\
$\t$ & 时间s \\
$\rpos$ & 物理径向距离m \\
$\R$ & 当前半径m \\
$\Rzero$ & 初始半径m \\
$\Llen$ & 长度m \\
$\rho$ & 题定经验密度，用于有效热容量 \\
$\cp$ & 题定比热J/(kg K) \\
$\kcond$ & 导热系数W/(m K) \\
$\Ddiff$ & 有效水分扩散系数m²/s \\
$\hconv$ & 对流换热系数W/(m² K) \\
$\hmass$ & 固相基准的有效供水系数m/s \\
$\Tinf$ & 烘房温度K \\
$\Yinf$ & 空气湿含量kg水/kg干空气；Boundary.C\_amb为兼容旧接口的存储名 \\
$\Ysat$ & 表面温度下饱和空气湿含量 \\
$\Yint$ & 反算的界面空气湿含量 \\
$\RH$ & 相对湿度0—1 \\
$\pair$ & 空气总压Pa \\
$\pvap$ & 水蒸气分压Pa \\
$\psat$ & 饱和蒸汽压Pa \\
$\epsratio$ & 水蒸气/干空气分子量比0.621945 \\
$\cpm$ & 每kg干空气的湿空气比热1006+1860Y∞，J/(kg干空气 K) \\
$\Le$ & Lewis数，热扩散率/水蒸气扩散率 \\
$\KM$ & 气膜湿含量驱动系数kg干空气/(m² s) \\
$\Ts$ & 药材表面温度 \\
$\Cs$ & 药材表面干基含水率 \\
$\rhod$ & 参考干物质密度kg/m³ \\
$\rhodzero$ & 初始参考干密度rho(C0,T0)/(1+C0) \\
$\jvap$ & 向外相变质量通量kg/(m² s)，负号为冷凝 \\
$\jsupply$ & 内部供水能力kg/(m² s) \\
$\jgas$ & 气膜排湿能力kg/(m² s) \\
$\Lvap$ & 纯水汽化潜热J/kg \\
$\wfrac$ & 湿基水质量分数C/(1+C) \\
$\xi$ & 材料坐标r/R \\
$\vmat$ & 材料径向速度r Rdot/R，m/s \\
$\Mdry$ & 干物质总质量kg \\
$\Mwater$ & 水总质量kg \\
$\Cmax$ & 全内部场最大含水率 \\
$\tstar$ & 全场达到0.15的连续阈值时刻s \\
$\dt$ & 时间步长s，必须区分内部与输出 \\
$\dr$ & 径向间距m \\
$\Nr$ & 内部径向节点数 \\
$\Vi$ & 无量纲环形控制体积权重 \\
$\FT,\,\FC$ & 材料域的面通量，见numerics.md \\
\bottomrule\end{longtable}
